%------------------------------------
% Short CV for manual, application-specific use.
%------------------------------------

%!TEX TS-program = xelatex
%!TEX encoding = UTF-8 Unicode

\documentclass[10pt, a4paper]{article}
\usepackage{fontspec}
\usepackage[autostyle=true]{csquotes}

% DOCUMENT LAYOUT
\usepackage{geometry}
\geometry{a4paper, textheight=27cm, left=2cm, right=2cm, top=2cm, bottom=2cm}
\setlength\parindent{0in}

% COLORS
\usepackage[dvipsnames]{xcolor}
\definecolor{Jet}{HTML}{2A2D31}
\definecolor{DarkCornflowerBlue}{HTML}{2C446F}
\definecolor{ViridianGreen}{HTML}{468F92}
\definecolor{SlateGray}{HTML}{6C7F93}

% FONTS
\defaultfontfeatures{Mapping=tex-text}
\setromanfont[
    Path = ../assets/fonts/,
    Extension = .otf,
    UprightFont = *-Regular,
    BoldFont = *-Bold,
    ItalicFont = *-Italic,
    BoldItalicFont = *-BoldItalic,
    Ligatures = Common,
    Numbers = OldStyle
]{ACaslonPro}

\newcommand{\amper}{{\itshape\addfontfeatures{Scale=.95}\&}}

% CUSTOM MACROS
\newlength{\datebox}
\setlength{\datebox}{2.8cm}
\newlength{\textbox}
\setlength{\textbox}{14.2cm}

\newcommand{\workitem}[4]{
    \noindent
    \parbox[t]{\datebox}{#1}%
    \parbox[t]{\textbox}{
        \textbf{\textcolor{ViridianGreen}{#2}}\\
        \textit{#3}\\
        \textcolor{Jet}{\small #4}
    }\par\vspace{0.4cm}
}

\newcommand{\eduitem}[4]{
    \noindent
    \parbox[t]{\datebox}{#1}%
    \parbox[t]{\textbox}{
        \textbf{\textcolor{ViridianGreen}{#2}} | \textit{#3}\\
        \textcolor{Jet}{\small #4}
    }\par\vspace{0.35cm}
}

\newcommand{\certitem}[3]{
    \noindent
    \parbox[t]{\datebox}{#1}%
    \parbox[t]{\textbox}{
        \textbf{\textcolor{SlateGray}{#2}} \textit{#3}
    }\par\vspace{0.25cm}
}

\newcommand{\skillitem}[2]{
    \noindent\textbf{\textcolor{SlateGray}{#1}}: #2\par\vspace{0.2cm}
}

\newcommand{\listitem}[3]{
    \noindent
    \parbox[t]{\datebox}{#1}%
    \parbox[t]{\textbox}{
        \textbf{\textcolor{ViridianGreen}{#2}}\\
        \textcolor{Jet}{\small #3}
    }\par\vspace{0.3cm}
}

% HEADINGS
\usepackage{sectsty}
\usepackage[normalem]{ulem}
\sectionfont{\rmfamily\mdseries\Large\textcolor{DarkCornflowerBlue}}
\subsectionfont{\rmfamily\mdseries\scshape\normalsize}
\subsubsectionfont{\rmfamily\bfseries\upshape\normalsize}

% PDF SETUP
\usepackage[xetex, bookmarks, colorlinks, breaklinks,
    pdftitle={Jorge Anais - Short Curriculum Vitae}, pdfauthor={Jorge Anais}]{hyperref}
\hypersetup{linkcolor=teal,citecolor=teal,filecolor=black,urlcolor=teal}

\usepackage{multicol}

% BIBLIOGRAPHY SETUP
\usepackage[
    backend=biber,
    style=numeric,
    sorting=ydnt,
    maxbibnames=99
]{biblatex}
\addbibresource{../data/publications.bib}

\begin{document}
{\LARGE \textcolor{DarkCornflowerBlue}{Jorge Anais Vilchez}} \\

\begin{minipage}{\linewidth}
\begin{multicols}{2}
Nationality: Chilean\\
Address: Santiago, Chile.\\
\vfill\eject
Email: \href{mailto:jrganais@gmail.com}{jrganais@gmail.com}\\
\textsc{url}: \href{http://www.jorgeanais.cl}{http://www.jorgeanais.cl}\\
\end{multicols}
\end{minipage}

\hrule
\vspace{0.3cm}

\begin{flushright}
  \begin{minipage}{15cm}
  \textcolor{Jet}{\small Astronomer and Data Scientist with expertise in deep learning, unsupervised clustering, and big data processing pipelines. Proven track record of designing scalable, data-driven solutions for complex environments, ranging from corporate mobility to wide-field astronomical surveys. Highly skilled in translating technical concepts for stakeholders and leading academic instruction in artificial intelligence.}
  \end{minipage}
\end{flushright}

\section*{Education}
\eduitem{2022--2027}{Doctorate in Astrophysics and Astroinformatics}{Universidad de Antofagasta}{PhD candidate. Thesis: \enquote{Characterization of Sagittarius dwarf galaxy at low Galactic latitudes}}
\eduitem{2018--2020}{Master in Astronomy}{Universidad de Antofagasta}{Thesis: \enquote{Search and characterization of star clusters in the far end of the Galactic bar}.}
\eduitem{2008--2014}{Licentiate in Astronomy}{Pontificia Universidad Católica de Chile}{Graduation activity: \enquote{A new astrometric software for the Swope Telescope}, developed at Las Campanas Observatory.}

\section*{Research \amper{} Professional Experience}
\workitem{2023--Present}{Research: Characterization of the Sagittarius dwarf spheroidal galaxy}{Universidad de Antofagasta}{Principal investigator of the project funded by the ANID National Doctoral Scholarship (2024-2026). Detailed characterization of the structure, stellar content, and dynamics of the Sagittarius dwarf galaxy in low Galactic latitude regions, combining photometric and spectroscopic data from multiple surveys.}
\workitem{2023--Present}{Assistant Professor}{Duoc UC, School of Informatics and Telecommunications}{Teaching Machine Learning and Deep Learning courses for computer engineering students.}
\workitem{04/25--10/25}{Early Career Visitor Programme}{ESO, Chile}{Researcher in the project led by Dr. Elisa R. Garro combining deep HST photometry and FLAMES+UVES spectroscopy to precisely determine ages and chemical compositions of Sagittarius dwarf galaxy globular clusters.}
\workitem{09/2024}{Research: Variable stars in the Sagittarius dwarf spheroidal galaxy}{NOIRLab}{Collaboration with NOIRlab researcher Kathy Vivas to detect and classify variable stars (focus on $\delta$ Scuti and RR Lyrae stars) in the central regions of Sagittarius galaxy using multiepoch DECam data.}
\workitem{2021--2022}{Data Scientist}{Grupo Kaufmann AWTO}{Design and implement machine learning models, perform exploratory data analysis, communicate insights and recommendations to stakeholders.}
\workitem{03/20--11/20}{Research: Search and characterization of star clusters}{CITEVA, Universidad de Antofagasta}{Development of a methodology based on unsupervised machine learning to search for young cluster candidates at the far end of the Galactic bar. Under the supervision of Dr. Sebastián Ramírez and Dra. Karla Peña Ramírez.}
\workitem{03/20--11/20}{Research Assistant: Spectroscopy models for planetary atmospheres}{CITEVA, UA}{Study on transmission spectroscopy \textit{in silico} models for Earth-like and Mini-Neptunes planetary atmospheres. Under the supervision of Dr. Jeremy Tregloan-Reed.}
\workitem{2015--2025}{Observer: 1M2H Project}{Las Campanas Observatory, UC Santa Cruz}{Photometric transients follow-up using Swope Telescope. Responsible for data acquisition, data quality assurance, and system troubleshooting.}
\workitem{01/16--06/16}{Research Assistant: Data reduction}{Instituto de Astronomía, UCN}{Data reduction of VLT FORS and Magellan MagE data. Under the supervision of Dr. Christian Moni Bidin.}
\workitem{2014--2015}{Observer: Carnegie Supernova Project II}{Las Campanas Observatory}{Observations of low-redshift supernovae using the Swope Telescope.}
\workitem{2014--2015}{Research Assistant: Sonification of medical images}{Instituto de Música, PUC}{Development of a sonification tool to assist radiologists in the detection of breast cancer.}

\section*{Observational Experience}
\listitem{2014--2025}{Swope 1m Telescope | Las Campanas Observatory}{More than 300 nights observed.}
\listitem{2018--2020}{Chakana 0.6m Telescope | UA Ckoirama Observatory}{10 nights observed in total.}
\listitem{2020}{Zeiss 1.23m Telescope | Calar Alto Astronomical Observatory}{8 nights observed in total (remote).}
\listitem{2013--2014}{PUC 50cm Telescope / PUCHEROS Spectrograph | PUC Santa Martina}{10 nights observed in total.}

\section*{Teaching \amper{} Course Development}
\workitem{2023--Present}{Instructor}{Duoc UC, School of Informatics and Telecommunications}{Teaching Machine Learning and Deep Learning courses for computer engineering students.}
\workitem{2010--2014}{Teaching Assistant}{Pontificia Universidad Católica de Chile}{Courses: Introduction to Physics for Biological Sciences, Static and Dynamics for Civil Engineers, Thermodynamics, Optics, Astronomical Instrumentation and Electricity and Magnetism.}

\subsection*{Courses Taught (Duoc UC)}
\begin{itemize}
    \item \textbf{Autumn 2026}: AVY1101 Big Data, AVY1102 Fundamentals of Machine Learning, BIY7121 Data Mining
    \item \textbf{Spring 2025}: MLY0100 Machine Learning
    \item \textbf{Autumn 2025}: DLY0100 Deep Learning, FMY0100 Fundamentals of Machine Learning
    \item \textbf{Spring 2024}: MLY0100 Machine Learning
    \item \textbf{Autumn 2024}: DLY0100 Deep Learning
    \item \textbf{Autumn 2023}: DLY0100 Deep Learning
    \item \textbf{Spring 2023}: MLY0100 Machine Learning
\end{itemize}

\subsection*{Courses Developed (Duoc UC)}
\begin{itemize}
    \item \textbf{2026}: MLT2202 Neural Networks and Deep Learning (Data Scientist Program online)
    \item \textbf{2025}: TLY1101 Machine Learning (Computer Engineering with AI specialization)
\end{itemize}

\section*{Certifications \amper{} Further Degrees}
\certitem{2026}{Data Governance, Regulation, and Ethics.}{Universidad de Chile}
\certitem{2025}{Deep Learning Specialization.}{Coursera - DeepLearning.AI}
\certitem{2025}{Evolutionary Algorithms (CC66M).}{Universidad de Chile}
\certitem{2025}{Transformers, use of agents, LangChain and RAG.}{Universidad de Chile}
\certitem{2024}{Visual Recognition with Deep Learning (CC66T).}{Universidad de Chile}
\certitem{2024}{Natural Language Processing (CC66Q).}{Universidad de Chile}
\certitem{2021}{Diplomature in Data Science.}{Universidad del Desarrollo}
\certitem{2018}{Diplomature in Software Development.}{Duoc UC}
\certitem{2016}{Diplomature in Astroengineering.}{Universidad de Antofagasta}

\section*{Achievements \amper{} Scholarships}
\begin{itemize}
    \item \textbf{2024--2026}: Full Scholarship for advanced human capital in Chile, Agencia Nacional de Investigación y Desarrollo de Chile.
    \item \textbf{2022--2023}: Academic Excellence Scholarship for PhD students, Universidad de Antofagasta.
    \item \textbf{2019}: Full scholarship to attend La Serena School for Data Science, AURA Observatory, Chile.
\end{itemize}

\section*{Conferences \amper{} Presentations}
\workitem{04/2025}{12th VVV Science Meeting: the VVVX synergy with next-generation Milky Way studies}{Universidade Federal de Santa Catarina, Brazil}{Oral: \textit{Unveiling the Sagittarius Dwarf Galaxy: Mapping Its Core and Streams at Low Galactic Latitudes}}
\workitem{07/2024}{Summer School in Statistics for Astronomers}{Center for Astrostatistics at The Pennsylvania State University}{}
\workitem{03/2024}{First LSST Latin American Meeting (LSST@LATAM)}{La Serena, Chile}{}
\workitem{03/2023}{XVII Latin American Regional IAU Meeting}{Montevideo, Uruguay}{Poster contribution: \textit{Unraveling the Sagittarius Dwarf Spheroidal Galaxy at Low Galactic latitudes}}
\workitem{03/2023}{First Latinamerican Virtual Observatory School}{Spanish Virtual Observatory (Online)}{}
\workitem{2021}{Pycon Chile 2021}{Online (\url{https://pycon.cl/2021/})}{Contributed talk: \textit{Search for star clusters using unsupervised machine learning}}
\workitem{2019}{LARIM 2019: XVI Latin American Regional IAU Meeting}{Antofagasta, Chile}{}
\workitem{2019}{Star Form Mapper 2019}{York, United Kingdom}{Contributed poster: \textit{Massive Open clusters in VVV data using unsupervised clustering algorithms}}
\workitem{2019}{La Serena School for Data Science: Applied Tools for Data-driven Sciences}{AURA, La Serena, Chile}{}

\section*{Skills \amper{} Languages}
\skillitem{Software}{Proficient in developing data analysis, reduction, and visualization using Python and its main associated modules (Astropy, Matplotlib, Numpy, Pandas, Plotly, PyTorch, Seaborn, SciPy). Astronomical software: IRAF, DS9, TOPCAT. Other technologies: Bash, Docker, Git, \LaTeX, OpenClaw, OpenCode, PostgreSQL, R. Github repo: \url{https://github.com/jorgeanais}}
\skillitem{Languages}{Spanish (native), English (professional competence).}

\section*{Publications}
\skillitem{ADS}{Full list on ADS: \url{https://ui.adsabs.harvard.edu/public-libraries/veXcb7VPROqU0r03axnXkg}}

\vspace{1cm}
\begin{center}
{\small \today\- \textbullet\- Santiago, Chile}\\
{\scriptsize Full Curriculum Vit\ae}
\end{center}

\end{document}
